%
% 53-homologie.tex -- Homologie von Ketten
%
% (c) 2026 Prof Dr Andreas Müller
%
\subsection{Homologie von Ketten
\label{buch:integration:homologie:subsection:homologie}}
Die Homotpieeigenschaft des Wegintegrals hat gezeigt, dass sich der
Wert eines Wegintegrals einer holomorphen Funktion nicht ändert, wenn
man den Weg stetig deformiert.
Dies wurde zum Beispiel in Abbildung~\ref{buch:integration:fig:residuum}
verwendet, um zu zeigen, dass der Weg $\gamma$ und der kleine Kreis
um den Punkt $z_0$ zu den gleichen Werten der Wegintegrale führen.
Diese Äquivalenz von Ketten, die zu gleichen Wegintegralen führen,
verdient einen eigenen Namen.

\begin{definition}[homolog]
Die Ketten $c_1$ und $c_2$ heissen \emph{homolog} in $U$, in Zeichen
$c_1\sim c_2$ oder $c_1\sim_{U}c_2$, wenn
\index{homolog}%
für jeden Punkt $a\in \mathbb{C}\setminus U$
\[
\operatorname{ind}_{c_1}(a)
=
\operatorname{ind}_{c_2}(a)
\]
gilt.
\end{definition}

Die Kette $\gamma_++\gamma_-$ ist also zu $\gamma$ homolog.
Für homologe Ketten ist nicht nur die Umlaufzahl gleich, sondern
jedes Integral.

\begin{satz}
Sind die Ketten $c_1$ und $c_2$ homolog, dann ist
\[
\int_{c_1} f(z)\,dz
=
\int_{c_2} f(z)\,dz.
\]
\end{satz}

Wir betrachten jetzt einen später besonders nützlichen Spezialfall.
\input{chapters/030-integration/fig/fig-merohomolog.tex}%
In Abbildung~\ref{buch:integration:homologie:fig:merohomolog}
ist ein Gebiet $U$ mit einer einfach geschlossenen Kurve
$\gamma$ dargestellt.
Das Komplement von $U$ im Inneren der Kurve $\gamma$ besteht
aus den isolierten Punkten $z_k$.
Jeder dieser Punkte kann so mit der Kurve $\gamma$ verbunden werden,
dass sich die Verbindungen nicht schneiden.
Ausserdem wird jeder von diesen Punkten von einem kleinen Kreis $\gamma_k$
im Gegenuhrzeigersinn umlaufen.
Aus $\gamma$, den im Uhrzeigersinn durchlaufenen Kreisen $-\gamma_k$
und den Verbindungen zu $\gamma$ lässt sich ein neuer Weg zusammensetzen,
der sich in $U$ zusammenziehen lässt.
Das Wegintegral über den gesamten Weg verschwindet daher.
Die Verbindungswege zur Kurve $\gamma$ werden dabei zweimal in
entgegengesetzter Richtung durchlaufen und tragen daher nicht zum
Wegintegral bei.
Das Wegintegral einer in $U$ holomorphen Funktion über $\gamma$
ist daher gleich der Summe der Wegintegrale der Funktion
über die Kreise $\gamma_k$.

\begin{satz}
\label{buch:integration:homologie:satz:pfadpunkte}
Ist $\gamma$ eine einfach geschlossene Kurve in $U$ und besteht das Komplement
von $U$ im Inneren von $\gamma$ nur aus endlich vielen isolierten Punkten
$z_k$, $1\le k\le n$, dann gibt es einen Radius $r$ derart, dass $\gamma$
homolog zu einem Zyklus 
\[
c
=
\sum_{k=1}^n \gamma_k,
\]
wobei jedes $\gamma_k$ ein Kreis vom Radius $r$ um $z_k$ ist.
\end{satz}

\begin{proof}
Es muss nur der Radius $r$ berechnet werden.
Nimmt man
\[
r
=
\frac13
\min \{ |z_k-z_l|\mid 1\le k,l\le n\},
\]
dann können sich die Kreise um die Punkte $z_k$ mit Radius $r$ nicht
schneiden.
Für jeden Punkt $z_k$ ist dann die Umlaufzahl 
\[
\operatorname{ind}_{\gamma_k}(z_l)
=
\delta_{kl}
=
\begin{cases}
1&\qquad\text{für $k=l$}\\
0&\qquad\text{für $k\ne l$.}
\end{cases}
\]
Für den ganzen Zyklus $C$ gilt dann
\[
\operatorname{ind}_c(z_l)
=
\sum_{k=1}^n
\delta_{kl}
=
1
=
\operatorname{ind}_\gamma(z_l).
\]
Dies zeigt, dass $\gamma$ und $c$ homolog sind.
\end{proof}
